% =============================================================================
% Additions for blocks 6-7 (Lectures 6-7, items 51-70) of the real-exercise list.
% Input right after preamble.tex by lista-reais-lecture-06.tex, lista-reais-lecture-07.tex and
% lista-reais-lectures-01-07.tex. Blocks 6-7 do not retype statements: each item \input's the
% fragment that Listas/make_statements.py generates from its practice list
% (pratica-lecture-0N-solutions/stmt/item-NN.tex), so the practice list stays the single source.
% Rerun make_statements.py after editing a practice list, then recompile.
% =============================================================================
\usepackage{longtable,calc}
\makeatletter
% Support macros that pandoc's LaTeX output expects (same as Listas/statement-hook.tex).
\providecommand{\tightlist}{\setlength{\itemsep}{0pt}\setlength{\parskip}{0pt}}
\@ifundefined{c@none}{\newcounter{none}}{}
\providecommand{\pandocbounded}[1]{\begingroup
  \setkeys{Gin}{width=0.75\linewidth,height=0.32\textheight,keepaspectratio}#1\endgroup}
\makeatother
\DeclareUnicodeCharacter{2212}{\ensuremath{-}}
\DeclareUnicodeCharacter{00B3}{\textsuperscript{3}}
% Fragments name images ../figures/x.png, relative to their solutions folder.
\graphicspath{{pratica-lecture-08-solutions/}}

% \exercicioreal[<diff>]{<short name>}{<level>}{<origin>}: header of a block 6-7 item. The
% fragment that follows prints the book source line itself, so the header gives the origin.
\newcommand{\exercicioreal}[4][]{\needspace{12\baselineskip}\refstepcounter{questao}\medskip\noindent
  \textbf{Exercise \thequestao.}\ifperlecture~{\small(item
  \the\numexpr\value{questao}+\blockoffset\relax\ of the combined list)}\fi~\textit{#2}%
  {\unskip\nobreak\hfil\penalty50\hskip1em\hbox{}\nobreak\hfill\mbox{\small[#3; difficulty #1/5]}}\par
  \noindent{\small\color{black!70}\textbf{Origin.} #4}\par\smallskip}
% \stmtfrag{<lecture NN>}{<item NN>}: the verbatim statement block from that practice list.
\newcommand{\stmtfrag}[2]{{\sloppy\input{pratica-lecture-#1-solutions/stmt/item-#2}}\par}

\renewcommand{\howtouse}{{\small\noindent\textbf{How to use this list.} Every exercise below is a
real exercise from one of the course books, reproduced verbatim under its source line (author,
book, exercise, printed page and, in brackets, the page your PDF viewer shows for the file in
\texttt{Livros/}, and whether the source gives the answer). Material the exercise refers to is
copied into the same block. For Lectures 6--7 the statements are printed from the practice lists
\texttt{Listas/pratica-lecture-0N.md}, where they were transcribed and checked against the book
pages; the ``Origin'' line names the chapter and the practice-list item. Muñoz-Garcia \&
Toro-Gonzalez print a worked answer right after each exercise; MWG answers are in the MWG
solutions manual; Tadelis and Gibbons give no answer in the book. There is one block per lecture
of the real schedule, ten exercises in rising difficulty, rated from 1 (apply a definition) to 5
(graduate level): three foundations, five exam level, two advanced. Detailed solutions:
\texttt{lista-reais-lecture-0N-solutions.pdf}.\par}}

\renewcommand{\tablenote}{{\footnotesize \ifperlecture \#: item number in this lecture's list,
then its number in the combined list (\texttt{lista-reais-lectures-01-07.pdf}) in brackets.
\fi Blocks 1--5 are those of \texttt{lista-reais-lectures-01-05.pdf} (marks there: $^\dagger$\,added
in the 30-item version, $^\ddagger$\,added in the 50-item version); blocks 6--7 were added on
2026-11-20 from the slides of Lectures 6 and 7. ``Diff.'': difficulty from 1 (apply a definition)
to 5 (graduate level). Pages: printed page, then the PDF viewer page in brackets (offsets in
\texttt{Livros/book-registry.json}; MWG is a two-pages-per-sheet scan). ``Answer'': \emph{book} =
worked answer printed after the exercise; \emph{manual} = MWG solutions manual in
\texttt{Livros/}; \emph{none} = no answer in the book. Books: Tadelis (2013), \emph{Game Theory:
An Introduction}; Gibbons (1992), \emph{Game Theory for Applied Economists}; Muñoz-Garcia \&
Toro-Gonzalez (2019), \emph{Strategy and Game Theory: Practice Exercises with Answers}, 2nd ed.;
Mas-Colell, Whinston \& Green (1995), \emph{Microeconomic Theory} (MWG); Cabral (2017),
\emph{Introduction to Industrial Organization}, 2nd ed.\par}}

\newcommand{\rowsvi}{%
L6 & \rownum{51} & Gibbons 3.1 & 169 [175] & Found. & Definitions: static Bayesian game, strategy, BNE & 1 & none \\
   & \rownum{52} & Tadelis 11.1 & 236 [255] & Found. & Disagreement as a Nash equilibrium of any bargaining game & 1 & none \\
   & \rownum{53} & Muñoz-Garcia Ch.~4 Ex.~1 & 154 [166] & Found. & Ultimatum with envy and guilt: SPE as a function of $\alpha$ & 2 & book \\
   & \rownum{54} & Gibbons 3.4 & 169--170 [175--176] & Exam & Nature picks the game: all pure BNE & 2 & none \\
   & \rownum{55} & Muñoz-Garcia Ch.~7 Ex.~1 & 318--319 [329--330] & Exam & Complete vs.\ incomplete information; Bayesian normal form & 3 & book \\
   & \rownum{56} & Tadelis 11.2 & 236 [255] & Exam & Hold-up: SPE vs.\ a better Nash equilibrium & 3 & none \\
   & \rownum{57} & Tadelis 12.3 & 266--267 [285--286] & Exam & Two-sided private types: armed conflict, all pure BNE & 3 & none \\
   & \rownum{58} & Tadelis 12.2 & 266 [285] & Exam & Cournot with private costs (the slides' homework 1) & 4 & none \\
   & \rownum{59} & Tadelis 12.9 & 269 [288] & Adv. & Continuous types; purification of a mixed equilibrium & 4 & none \\
   & \rownum{60} & Tadelis 11.7 & 237 [256] & Adv. & Baron--Ferejohn with $T = 1, 2$ (the slides' homework 2) & 5 & none \\}
\newcommand{\rowsvii}{%
L7 & \rownum{61} & Tadelis 15.1 & 314 [333] & Found. & BNE vs.\ PBE in a small tree & 1 & none \\
   & \rownum{62} & Muñoz-Garcia Ch.~9 Ex.~1 & 376--377 [385--386] & Found. & Separating and pooling PBE, two types & 2 & book \\
   & \rownum{63} & Tadelis 15.3 & 315--316 [334--335] & Found. & Trust with player types; BNE that are not PBE & 2 & none \\
   & \rownum{64} & Tadelis 12.5 & 267--268 [286--287] & Exam & Selling a mine of unknown quality: adverse selection & 3 & none \\
   & \rownum{65} & Muñoz-Garcia Ch.~9 Ex.~2 & 385--386 [394--395] & Exam & Job-market signaling: four candidate profiles & 3 & book \\
   & \rownum{66} & Tadelis 12.6 & 268 [287] & Exam & Adverse selection in hiring; a policy that raises surplus & 3 & none \\
   & \rownum{67} & Tadelis 15.5 & 316--317 [335--336] & Exam & Entry with a tough or weak entrant: BNE and PBE & 4 & none \\
   & \rownum{68} & Tadelis 16.4 & 336--337 [355--356] & Exam & Productive education as a signal: BNE vs.\ PBE & 4 & none \\
   & \rownum{69} & MWG 13.C.5 & 475 [248] & Adv. & Regulated price: can advertising separate qualities? & 4 & manual \\
   & \rownum{70} & Tadelis 16.8 & 338 [357] & Adv. & Beer or quiche: pooling PBE and the intuitive criterion & 5 & none \\}

\newcommand{\pathvi}{\textbf{Lecture 6}: from the definitions of a Bayesian game (1) through
disagreement, hold-up and type-contingent strategies to Tadelis 11.7 (5), majority-rule
bargaining with a closed rule over one and two periods, compared with Rubinstein and with the
infinite-horizon Baron--Ferejohn solution.}
\newcommand{\pathvii}{\textbf{Lecture 7}: from separating BNE and PBE in a small tree (1) through
lemons-style unraveling and Spence's education game to Tadelis 16.8 (5), where every PBE is pooling
and the intuitive criterion must pick out the one with reasonable off-path beliefs.}
